% Part: Philosophical Introduction
% no need to use \olchapter and \OLEndChapterHook

\documentclass[../../../include/open-logic-chapter]{subfiles}

\begin{document}

\olpart{phi}{Philosophical Introduction}

\chapter{Three Puzzles Concerning Identity, Necessity and Existence}

This chapter introduces some philosophical puzzles concerning 
identity, necessity and existence. These puzzles have been
central to contemporary developments in logic, metaphysics, 
philosophy of mind and philosophy of language. 
They have spurred major developments in logic in the XXth century,
notably Quantified Modal Logic, but also Free Logic and Hyperintensional
Logics. These new logics have in turn improved our understanding 
of the puzzles. The new logics have also applications beyond the puzzles.

Here we introduce our puzzles without using any logical notation. 
You may notice that it is sometimes difficult to parse the claims
involved. That is intentional. Later in this course you'll be able 
to recast these puzzles using the language of quantified modal logic. 
Hopefully, you'll then experience how logical tools can shaprent philosophical thought.   

\section{Identity, Necessity and Existence}

Our three puzzles are the Puzzle of Necessary Existence,
the Puzzle of Necessary Identity, and Frege's Puzzle (which we 
might call the Puzzle of Unknown Identities).

All the puzzles involve the notion of \emph{identity}. What do we mean
by that? Being one and the same thing (also called \emph{numerical
identity}). For instance, Mount McKinley and Denali are one and the
same thing: the Alaskan mountain called `Mount McKinley' and the
Alaskan mountain called `Denali' are one and the same. `Identity' in
that sense is not something shared with others like being a Londoner,
or two cars being of the same make and model (what is sometimes called
\emph{qualitative identity}). It's a matter of being one and the same
thing. Thus two `identical' twins are not identical in the
philosophical sense: they are \emph{not} one and the same person.

Two of the puzzles invovle the notion of what is \emph{necessarily so} 
(\emph{metaphysical} necessity). What does that mean? We understand this notion in the 
following way: something is necessarily so iff it couldn't in any sense 
have been otherwise. For instance, it couldn't at all have been the 
case that there were cats and dogs without it being the case 
that there are cats. Therefore, the following is necessary: if there are
cats and dogs there are cats.

One puzzle also involve the notion of \emph{existence}. The moon exists; you exist;  so do some cats, some chairs, London, rivers, etc. Perhaps also some 
more esoteric things like viruses, atoms, and black holes. Unicorns, on 
the other hand, don't exist. But what is it to exist, more generally? 
One common idea is that \emph{to exist is to be 
something}. London is something, so it exists. Nothing is a unicorn, 
so unicorns don't exist. Another idea is that \emph{what exist is 
what there is}. There is a city that is London, so London exists. 
There are no unicorns, so unicorns don't exist. The two ideas aren't
necesssarily distinct: in fact, in first-order logic, we translate them 
in the same manner. 

\section{The Puzzle of Necessary Existence}

There is an argument from plausible principles to the conclusion that existence is necessary. 

The idea of something existing necessarily is often associated with
God. Many philosophers, theist and atheist alike, agree that if there
is in fact a God, then God is something that could not have failed to
exist. If the argument we're about to discuss is correct, however,
necessary existence isn't the prerogative of God: \emph{you} also
couldn't have failed to exist. That's a puzzling conclusion.

We first spell out a specific instance of the argument:

\begin{enumerate}
\item Necessarily, King Charles is King Charles.
\item Necessarily, if some thing is itself, then that thing is
something.\footnote{Here `thing' here is used in the most general
sense, aking to 'entity'. In that sense `thing' doesn't contrast with
animals or people.}
\item Necessarily, if some thing is something, then that thing exists. 
\item Necessarily, King Charles exists. (By 1-3 using \emph{Distribution})
\end{enumerate}

Where \emph{Distribution} (of Necessity) is the following principle:

\begin{quote}
\emph{Distribution}. If necessarily, if something is so then something is so, then: if necessarily, the former is so, then necessarily, the latter is so.
\end{quote}

For instance: if necessarily, if King Charles is a human then King
Charles is mortal, then: if necessarily, King Charles is human, then
necessarily, King Charles is mortal. Applying \emph{Distribution} to
the premises above, we can derive the conclusion:

\begin{itemize}
\item Necessarily, King Charles is himself. (From 1)
\item If necessarily, King Charles is himself, then necessarily, King
Charles is something. (From 2, \emph{Distribution})
\item If necessarily, King Charles is something, then necessarily,
King Charles exists. (From 3, \emph{Distribution})
\item Therefore, necessarily, King Charles exists.
\end{itemize}

The conclusion should strike you as obviously false. If it doesn't,
you are probably reading ``necessarily'' in its \emph{epistemic}, not
metaphysical sense: `given what we know, it cannot be otherwise
than...'. (We know that King Charles is alive, so given what we know
it cannot be otherwise than King Charles exist.) That is not the sense
of ``necessarily'' we're interested in here. We're interested in the
metaphysical sense: `it could not have been otherwise than'. To get
the intended sense of the argument, it may help you to consider a
version that makes this explicit:

\begin{enumerate}
\item It could not have happened that King Charles is not King Charles.
\item It could not have happened that some thing is itself and 
that thing is nothing.
\item It could not have happened that some thing is something and that
thing doesn't exist. 
\item Therefore, it couldn't have happened that King Charles doesn't
exist. 
\end{enumerate}

The conclusion seems obviously false. King Charles's parents could
have never met and King Charles never born. If so, it would seem, King
Charles wouldn't have existed. 

If you think that King Charles could have failed to exist, then you
reject the argument's conclusion. Therefore, you have to reject one of
its premises or the \emph{Distribution} principle. Otherwise, you have
to accept that not only King Charles exists, but it couldn't have
failed to exist. Which is it?

The argument for King Charles's necessary existence can obviously be
replicated for any existing thing, including you and me. We can state
it in a more general form:

\begin{enumerate}
    \item  Everything is necessarily itself. (\emph{Necessity of
    Self-Identity})
    \item Necessarily, if some thing is itself, then that thing is
    something. (\emph{Necessitated Existential Generalization})
    \item Necessarily, if some thing is something then that thing
    exists. (\emph{Necessitated Quantified Definition of Existence})
    \item Therefore, everything necessarily exists. (\emph{Necessity
    of Existence})
\end{enumerate}

In the recent literature, philosophers who endorse the necessity of
existence have been called \emph{Necessitists} and those who reject\
it \emph{Contingentist}.

Another reason why the argument is puzzling is that it infers a
necessity from a fact. In general, from the mere observation that
something is \emph{in fact} the case we can't infer that it
\emph{could not have been otherwise}. Hume famously thought this
couldn't be done: from the observation that one thing followed another
we can't infer that it had to follow the other. Yet if this argument
is right, Hume's line can be crossed: from the mere observation that
something exists, we can infer that it couldn't have failed to exist.
Therefore, if the argument is right, some truths are both \emph{a
posteriori} (discovered by  experience) and yet known to be
\emph{necessary} (such that they couldn't be otherwise). 

\section{The Puzzle of Necessary Identity}

There is an argument from plausible premises that identity is
necessary. This puzzling as it would seem that some identities are
contingent. 

Let us consider a specific instance of the argument first. Aurélie
Dupin wrote under the nom de plume George Sand. Hence George Sand and
Aurélie Dupin are one and the same person.

\begin{enumerate}
\item Necessarily, Aurélie Dupin is Aurélie Dupin.
\item If some thing is some thing, then everything that is true of the
former is true of the latter.\footnote{Here `thing' here is used in
the most general sense, aking to 'entity'. In that sense `thing'
doesn't contrast with animals or people.}
\item If George Sand is Aurélie Dupin, then necessarily, George Sand is Aurélie Dupin. (1,2 and \emph{De Re Modality Conversion})
\end{enumerate}

Where \emph{De Re Modality Conversion} is the following principle:

\begin{itemize}
\item \emph{De Re Modality Conversion} (DRMC). Necessarily some thing is some way if and
only if it is true of that thing that necessarily, it is in that way. 
\end{itemize}

The conclusion follows from the premises and the principle as follows:

\begin{itemize}
\item It is true of Aurélie Dupin that necessarily, she is Aurélie
Dupin. (by 1, DRMC)
\item If George Sand is Aurélie Dupin, then: if it is true of Aurélie 
Dupin that necessarily, she is Aurélie Dupin, then it is true of 
George Sand that necessarily, she is Aurélie Dupin.
Dupin. (by 2)
\item So, if George Sand is Aurélie Dupin, then: it is true of George Sand that necessarily, she is Aurélie Dupin. (by the two previous claims)
\item If it is true of George Sand that necessarily, she is Aurélie Dupin, then it is necessary that George Sand is Aurélie Dupin. (DRMC)
\item So, if George Sand is Aurélie Dupin, then it is necessary that George Sand is Aurélie Dupin. (by the two previous claims)
\end{itemize}

The conclusion should strike you as suprising at least, perhaps even
false. If it doesn't, you may be reading ``if ..., then necessarily,
..." in the sense of: ``from the fact or assumption that ... it
necessarily follows that ...'' (what Medieval logicians call the
\emph{necessity of consequence}). Understood in that way the
conclusion is not surprising but trivially true: from the fact that
George Sand is Aurélie Dupin it obviously follows that George Sand is
Aurélie Dupin. But the intended sense of the conclusion here is that
if George Sand is Aurélie Dupin, that that fact necessary (what
Medieval logicians call the \emph{necessity of the consequent}). That
is not obvious or trivial: for comparison, from the fact that I won a
lottery prize it doesn't follow that it was necessary that I would win
it. To get the intended sense of the argument, it may be useful to
rephrase it by making explicit the intended sense of `necessarily',
namely that something could not have been otherwise. That is:

\begin{enumerate}
\item It could not have been that Aurélie Dupin is not Aurélie Dupin.
\item If some thing is some thing, then everything that is true of the
former is true of the latter.\footnote{Here `thing' here is used in
the most general sense, aking to 'entity'. In that sense `thing'
doesn't contrast with animals or people.}
\item If George Sand is Aurélie Dupin, then it could not have been that
George Sand is not Aurélie Dupin. 
\end{enumerate}

The conclusion might seem false. Put yourself in the shoes of a reader
at the time. They might be surprised to learn that George Sand is 
Aurélie Dupin. They might have thought that George Sand could have 
been somebody else. 

The argument for the necessary identity of Aurélie Dupin and George 
Sand can obviously be replicated for any identity. Some classic examples: the Morning Star (Phosphorus) and the Evening Star (Hesperus), which turn
out to be Venus; Clark Kent and Superman; Mark Twain and Samuel Clemens. 
In each case you might seem that the identity is contingent. But if 
the argument is correct, that's false: these things couldn't have 
failed to be one and the same.

If you think these identities are contingent, you need to reject the
conclusion, and hence one of the premises or the \emph{De Re Modality
Conversion Principle} principle. Which is it? Otherwise, you need to
accept the conclusion: these things that are in fact identical
couldn't have failed to be identical. 

We can state the argument in a more general form:

\begin{itemize}
\item Everything is necessarily itself. \emph{Necessity of self-identity.} 
\item If something is one and the same thing as something, whatever is true of the former is true of the latter. \emph{Leibniz's law, a.k.a. the Indescernibility of Identicals}
\item Therefore, if something is one and the same thing as something, it is necessary that the former is one and the same thing as the latter. \emph{Necessity of identity}. 
\end{itemize}

Another reason why the argument is puzzling is that it infers a
necessity from a fact. In general, from the mere observation that
something is \emph{in fact} the case we can't infer that it
\emph{could not have been otherwise}. Hume famously thought this
couldn't be done: from the observation that one thing followed another
we can't infer that it had to follow the other. Yet if this argument
is right, Hume's line can be crossed: from the mere observation that
some thing is one and the same as some thing, we can infer that the 
former couldn't have failed to be the latter.
Therefore, if the argument is right, some truths are both \emph{a
posteriori} (discovered by  experience) and yet known to be
\emph{necessary} (such that they couldn't be otherwise). 

\section{Frege's puzzle}

There is an argument from plausible premises that all true identities
are known. This is puzzling as it seems obvious that some identities
are unknown. 

Consider the following story (accurate, for all I know, but I haven't
checked). At the times of Homer, Greeks had a name for the first
'star' (in fact, planet) to be visible at night, Hesperus. They also
had a name for the last 'star' to be visible in the morning,
Phosphorus. Only centuries later, Greek astronomers realized that
these where one and the same `star' (the planet we call Venus). If so,
both of the following claim seem true:

\begin{itemize}
    \item Phosphorus is (the same thing as) Hesperus.
    \item Homer did not know that Phosphorus is Hesperus.
\end{itemize}

However, there is an argument from that leads to the opposite
conclusion.

\begin{enumerate}
\item Homer knew that Hesperus is Hesperus. 
\item If some thing is some thing, then everything that is true of the
former is true of the latter.\footnote{Here `thing' here is used in
the most general sense, aking to 'entity'. In that sense `thing'
doesn't contrast with animals or people.}
\item If Phosphorus is Hesperus, then Homer knew that Phosphorus is Hesperus. (1,2 and \emph{De Re Knoweldge Conversion})
\end{enumerate}

Where \emph{De Re Knowledge Conversion} is the following principle:

\begin{itemize}
\item \emph{De Re Knowledge Conversion} (DRKC). One knows that some thing is some way if and
only if it is true of that thing that one knows it to be that way. 
\end{itemize}

The conclusion follows from the premises and the principle as follows:

\begin{itemize}
\item It is true of Hesperus that Homer knows it to be Hesperus. (by 1, DRKC)
\item If Phosphorus is Hesperus, then: if it is true of Hesperus that Homer knows it to be Hesperus, then it is true of Phosphorus that Homer knows it to be Hesperus. (by 2)
\item So, if Phosphorus is Hesperus, then: it is true of Phosphorus that Homer knows it to be Hesperus. (by the two previous claims)
\item If true of Phosphorus that Homer knows it to be Hesperus, then Homer knows that Phosphorus is Hesperus. (DRMC)
\item So, if Phosphorus is Hesperus, then Homer knows that Phosphorus is Hesperus. (by the two previous claims)
\end{itemize}

The conclusion seems obviously false. Homer didn't know that Phosphorus
is Hesperus.

The argument can be obviously generalized to any true identity.

If you reject the conclusion, you must reject one of the premises
or the De Re Knoweldge Conversion. Which is is? Otherwise you 
have to accept the knowledge, and say that no true identity is unknown.

Kripke's `A Puzzle About Belief' poses a particularly sharp version Of
the problem. In his story, Pierre grows up in France and learns about 
London through photos and descriptions in French. Based on this, 
he sincerly makes claims such as:

\begin{itemize}
\item Londres est jolie.
\end{itemize}

Which translates as `London is pretty'. Later he moves to work in 
the UK, with little awareness of where he is, learns English with
the locals, and ends up living in a ugly industrial part of London.
He sincerely utters things like:

\begin{itemize}
\item London is not pretty.
\end{itemize}

He is not aware that ``London'' and ``Londres'' name the same city.
The opinions he expressed with the word ``Londres'' haven't changed;
he would still assent to the sentence ``Londres est jolie''. 

The puzzle is: \emph{does Pierre believe that London is pretty?} If we only  considered the part of the story, we would say that Pierre believe 
that London is pretty. He learned about London through books, in French,
he calls it ``Londres'', but that's a name for London. When he sincerely
utters ``Londres is pretty'', we should say that:

\begin{itemize}
\item Pierre believes that London is pretty. 
\end{itemize}

On the other hand, if we look at the second part of the story, he's
learned about London by living in it; he calls it `London', which 
is its English name indeed. When he sincerely assert that `London 
is not pretty', we should say at least that:

\begin{itemize}
\item Pierre believes that London is not pretty. 
\end{itemize}

If we say both things, we're ascribing Pierre \emph{contradictory beliefs}: he believes one thing, and he believes its negation. That
is a first puzzle, because having contradictory beliefs is usually
irrational, but Pierre doesn't seem irrational here---merely poorly
informed.

But that's not all. Based on the second part of the story alone, because 
Pierre sincerly rejects `London is pretty', we should not only say 
that Pierre believes that London is not pretty, but also that he \emph{does not} have the belief that London is pretty:

\begin{itemize}
\item Pierre does not believe that London is pretty. 
\end{itemize}

But would mean a contradition not just in Pierre's mind, but in our 
own theory of Pierre's beliefs! 

There are many things we can all agree on about Pierre's beliefs:

\begin{itemize}
    \item He believes that there is city called `Londres' that is pretty.
    \item He believes that there is a city called `London' that is ugly.
    \item He believes that `London' and `Londres' are names of different cities.
    \item He does not believe that `London' and `Londres' name the same city.
    \item He does not believe that there is an ugly city called `Londres'.
    \item He does not believe that there is a pretty city called `London'. 
\end{itemize}

But as Kripke insists, these does not settle the question he raises.
If it makes sense talk of someone believing that London is pretty, 
it must make sense to ask whether Pierre does believe it or not. 
Which is it?

We can relate Kripke's Pierre case with our argument. Using the 
word `Londres' ourselves, an argument like
the one above can be used to show that if Pierre believes that Londres
is pretty, he believes that Londres is pretty. 

\begin{enumerate}
\item Pierre believes that London is pretty. 
\item If Homer believes that London is pretty, then it is true of London 
that Homer believes it to be pretty. (DRKC)
\item So, it is true that of London that Pierre believes it to be pretty. (1, 2)
\item If some thing is some thing, then everything that is true of the
former is true of the latter.
\item Londres is London.
\item So, it is true of Londres that Pierre believes it to be pretty. (3,4,5)
\item If it is true of Londres that Pierre believes it to be pretty, then Pierre believes Londres to be pretty. (DRKC)
\item So, Pierre believes that Londres is pretty.
\end{enumerate}

The conclusion seems false, as Pierre insists that the city he calls
Londres isn't pretty. If you reject the conclusion, you must reject
one of the premises, or the principle of \emph{De Re Knoweldge Conversion}.
Which is it? Otherwise, if you accept it, you must accept that Pierre
does believe that London is pretty. (And also, by parallel reasoning,
that he does believe that London is ugly.)

The argument behind Frege's puzzle is formally like the argument 
for the necessity of identity. If you think that the argument for 
the necessity of identity is correct, but the argument for 
the knowledge of identities is not, you should consider why
they differ.

The puzzle of unkonw identities was first raised by Frege. He raised
it in slightly different terms. He compares the sentences `Hesperus is
Hesperus' and 'Hesperus is Phosphorus'. He points out their difference
in informativeness: the former isn't informative (anyone would have
figured it out that it is true), the second is informative (it took
careful astronomic observations to establish that it was true).
However, the two sentences seem to have the same meaning (to `say' the
same thing): `Hesperus' denotes a certain planet, and `Phosphorus'
denotes the same planet. Yet their different informativeness must be
due to different meaning (to them `saying' different things). 

%\olimport[story]{story}
%\olimport[ordinals]{ordinals}

\end{document}